% GENERATED FILE. Edit canonical lessons, then run npm run generate:books.
% canonical-source-hash: fnv1a64:b63d71c537296fb8
% canonical-lessons: KA-C155-ittichege, KA-C155-igagale, KA-C155-kaleda, KA-C155-omme, KA-C155-kale

\chapter{Recently, Already, Last, Once, To Spend (Time)}
\label{ch:ka-recently-already-last-once-to-spend-time}
\hlchaptermodality{155}

\begin{chapteropening}
\textbf{By the end of this chapter:} \emph{I can say what happened, and when: recently, already, last year.}

Complete the last lesson of chapter 155: I can say what happened, and when: recently, already, last year.
\end{chapteropening}

\section[\texorpdfstring{\kn{ಇತ್ತೀಚೆಗೆ} (ittīcege)}{ಇತ್ತೀಚೆಗೆ (ittīcege)}]{\kn{ಇತ್ತೀಚೆಗೆ} (ittīcege) — recently}

\label{lesson:KA-C155-ittichege}



\begin{quote}
Before the new one: say the Kannada for anyone, then the Kannada for anything.
\end{quote}

\subsection*{You'll want to know: \kn{ಇತ್ತೀಚೆಗೆ}}
\textbf{\kn{ಇತ್ತೀಚೆಗೆ}} — \emph{ittīcege} — \textquotedblleft{}recently\textquotedblright{}. Kannada's past for \textquotedblleft{}I\textquotedblright{} ends in \textbf{-\kn{ಎ}} (\emph{-e}) after the past stem: \textbf{\kn{ನಾನು} \kn{ಇತ್ತೀಚೆಗೆ} \kn{ಒಂದು} \kn{ಪುಸ್ತಕ} \kn{ಓದಿದೆ}} — \emph{nānu ittīcege ondu pustaka ōdide} — \textquotedblleft{}I recently read a book.\textquotedblright{}

\begin{grammarlens}[title={What you've built}]
One more word for this chapter.
\end{grammarlens}

\subsection*{Your turn}
\begin{itemize}
\raggedright
  \item \emph{Say it:} \emph{ittīcege}
  \item \emph{Say it:} \emph{ittīcege}, once more
  \item \emph{Say it:} say \emph{ittīcege}
  \item \emph{Recall:} say the Kannada for anyone, then the Kannada for anything, then say \emph{ittīcege} again
\end{itemize}

\begin{tcolorbox}[breakable,colback=teal!4,colframe=teal!35!black,title={Before you move on}]
What does \kn{ಇತ್ತೀಚೆಗೆ} mean? (\textquotedblleft{}Recently\textquotedblright{}.) Say it once more.
\end{tcolorbox}
\section[\texorpdfstring{\kn{ಈಗಾಗಲೇ} (īgāgalē)}{ಈಗಾಗಲೇ (īgāgalē)}]{\kn{ಈಗಾಗಲೇ} (īgāgalē) — already}

\label{lesson:KA-C155-igagale}



\begin{quote}
Before the new one: say the Kannada for anything, then the Kannada for recently.
\end{quote}

\subsection*{You'll want to know: \kn{ಈಗಾಗಲೇ}}
\textbf{\kn{ಈಗಾಗಲೇ}} — \emph{īgāgalē} — \textquotedblleft{}already\textquotedblright{}. \textbf{\kn{ನಾನು} \kn{ಈಗಾಗಲೇ} \kn{ಊಟ} \kn{ಮಾಡಿದೆ}} — \emph{nānu īgāgalē ūṭa māḍide} — \textquotedblleft{}I have already eaten.\textquotedblright{}

\begin{grammarlens}[title={What you've built}]
One more word for this chapter.
\end{grammarlens}

\subsection*{Your turn}
\begin{itemize}
\raggedright
  \item \emph{Say it:} \emph{īgāgalē}
  \item \emph{Say it:} \emph{īgāgalē}, once more
  \item \emph{Say it:} say \emph{īgāgalē}
  \item \emph{Recall:} say the Kannada for anything, then the Kannada for recently, then say \emph{īgāgalē} again
\end{itemize}

\begin{tcolorbox}[breakable,colback=teal!4,colframe=teal!35!black,title={Before you move on}]
What does \kn{ಈಗಾಗಲೇ} mean? (\textquotedblleft{}Already\textquotedblright{}.) Say it once more.
\end{tcolorbox}
\section[\texorpdfstring{\kn{ಕಳೆದ} (kaḷeda)}{ಕಳೆದ (kaḷeda)}]{\kn{ಕಳೆದ} (kaḷeda) — last, previous}

\label{lesson:KA-C155-kaleda}



\begin{quote}
Before the new one: say the Kannada for recently, then the Kannada for already.
\end{quote}

\subsection*{You'll want to know: \kn{ಕಳೆದ}}
\textbf{\kn{ಕಳೆದ}} — \emph{kaḷeda} — \textquotedblleft{}last, previous\textquotedblright{}. \textbf{\kn{ಕಳೆದ} \kn{ವರ್ಷ} \kn{ನಾನು} \kn{ಮೈಸೂರಿಗೆ} \kn{ಹೋದೆ}} — \emph{kaḷeda varṣa nānu maisūrige hōde} — \textquotedblleft{}Last year I went to Mysuru.\textquotedblright{}

\begin{grammarlens}[title={What you've built}]
One more word for this chapter.
\end{grammarlens}

\subsection*{Your turn}
\begin{itemize}
\raggedright
  \item \emph{Say it:} \emph{kaḷeda}
  \item \emph{Say it:} \emph{kaḷeda}, once more
  \item \emph{Say it:} say \emph{kaḷeda}
  \item \emph{Recall:} say the Kannada for recently, then the Kannada for already, then say \emph{kaḷeda} again
\end{itemize}

\begin{tcolorbox}[breakable,colback=teal!4,colframe=teal!35!black,title={Before you move on}]
What does \kn{ಕಳೆದ} mean? (\textquotedblleft{}Last, previous\textquotedblright{}.) Say it once more.
\end{tcolorbox}
\section[\texorpdfstring{\kn{ಒಮ್ಮೆ} (omme)}{ಒಮ್ಮೆ (omme)}]{\kn{ಒಮ್ಮೆ} (omme) — once, one time}

\label{lesson:KA-C155-omme}



\begin{quote}
Before the new one: say the Kannada for already, then the Kannada for last.
\end{quote}

\subsection*{You'll want to know: \kn{ಒಮ್ಮೆ}}
\textbf{\kn{ಒಮ್ಮೆ}} — \emph{omme} — \textquotedblleft{}once, one time\textquotedblright{}. \textbf{\kn{ನಾನು} \kn{ಒಮ್ಮೆ} \kn{ದೆಹಲಿಗೆ} \kn{ಹೋಗಿದ್ದೆ}} — \emph{nānu omme dehalige hōgidde} — \textquotedblleft{}I went to Delhi once.\textquotedblright{}

\begin{grammarlens}[title={What you've built}]
One more word for this chapter.
\end{grammarlens}

\subsection*{Your turn}
\begin{itemize}
\raggedright
  \item \emph{Say it:} \emph{omme}
  \item \emph{Say it:} \emph{omme}, once more
  \item \emph{Say it:} say \emph{omme}
  \item \emph{Recall:} say the Kannada for already, then the Kannada for last, then say \emph{omme} again
\end{itemize}

\begin{tcolorbox}[breakable,colback=teal!4,colframe=teal!35!black,title={Before you move on}]
What does \kn{ಒಮ್ಮೆ} mean? (\textquotedblleft{}Once, one time\textquotedblright{}.) Say it once more.
\end{tcolorbox}
\section[\texorpdfstring{\kn{ಕಳೆ} (kaḷe)}{ಕಳೆ (kaḷe)}]{\kn{ಕಳೆ} (kaḷe) — to spend (time)}

\label{lesson:KA-C155-kale}



\begin{quote}
Before the new one: say the Kannada for last, then the Kannada for once.
\end{quote}

\subsection*{You'll want to know: \kn{ಕಳೆ}}
\textbf{\kn{ಕಳೆ}} — \emph{kaḷe} — \textquotedblleft{}to spend (time)\textquotedblright{}. Said of time: \textbf{\kn{ನಾನು} \kn{ರಜೆಯನ್ನು} \kn{ಹಳ್ಳಿಯಲ್ಲಿ} \kn{ಕಳೆದೆ}} — \emph{nānu rajeyannu haḷḷiyalli kaḷede} — \textquotedblleft{}I spent the holiday in the village.\textquotedblright{}

\begin{grammarlens}[title={What you've built}]
One more word for this chapter.
\end{grammarlens}

\subsection*{Your turn}
\begin{itemize}
\raggedright
  \item \emph{Say it:} \emph{kaḷe}
  \item \emph{Say it:} \emph{kaḷe}, once more
  \item \emph{Say it:} say \emph{kaḷe}
  \item \emph{Recall:} say the Kannada for last, then the Kannada for once, then say \emph{kaḷe} again
\end{itemize}

\begin{tcolorbox}[breakable,colback=teal!4,colframe=teal!35!black,title={Before you move on}]
What does \kn{ಕಳೆ} mean? (\textquotedblleft{}To spend (time)\textquotedblright{}.) Say it once more.
\end{tcolorbox}
